\begin{afterword}
\summaryhead

The post adds a new kind to the author's list of ``improper'' beliefs, meaning beliefs that do not guide what a person expects to experience. The new kind is belief worn as a sign of group membership. Following a metaphor of Robin Hanson's, the author calls it ``belief as attire.''

His example is the September 11 attacks. He says the hijackers surely saw themselves as heroes, yet in the United States one cannot say so, even to describe their view. Calling them cowards is American attire; calling them brave is the Enemy's.

He then suggests that attire explains why people can be passionate about beliefs that no longer guide their expectations. Merely believing that one ought to believe, he suspects, produces little real feeling. Belonging to a group produces a great deal, and the group's beliefs are spoken with that passion. A footnote calls the Republican--Democrat contest a ``swindle'' built on this.

\responsehead

The post names an old observation, that people hold some beliefs as badges, and gives it no method. Its one test for telling attire from conviction is a joke: a belief is Enemy attire ``because the Enemy talks about it.'' Without a test the label can be applied to any belief a reader dislikes, and the post's examples show where it will be applied.

Look at whose beliefs are examined. A believer praying for a sick baby. The hijackers and their ``Islamic Way.'' Religious people whose passion is ``not the same fire they had as a child.'' Americans in an Alabama bar. Partisans of either party, in the footnote. The reader and the author appear in the post only as the ones who see through all of these. What the post offers its reader is membership in a group whose beliefs are never attire.

Apply the label to the post itself. It has a vocabulary that marks the in-group (proper and improper belief, anticipation-controller, belief in belief), a shared enemy (the unreflective believer), and a footnote whose stance, above both parties, costs nothing to adopt and signals that the reader is not one of the swindled. By the post's own description, that is attire. Nowhere in the post does the author ask whether any belief of his own, or of his readers, is worn this way. This is an inference from what the post omits, but the omission is complete.

The central example is also false as written. The post says that in America one ``cannot'' call the hijackers brave. Within two weeks of the attacks Bill Maher said on television that staying in the airplane was ``not cowardly,'' and Susan Sontag wrote in \textsc{The New Yorker} that ``they were not cowards.'' Both paid a price for it, which is the true and weaker claim. The post does not cite them. It replaces the record with an imagined bar in Alabama, a stereotype that lets the reader feel superior to the people it describes.

Its claim about passion, finally, rests on an ``impression'' of religious people that no observation could overturn. It fails the rent test the author had published five days before.

In short: a label with no test, demonstrated on a false claim about American speech, applied to every group except the author's, and itself worn as attire.
\end{afterword}
